\section{Link del repositorio del trabajo}
\textbf{\textit{Puede ser Github, Gitlab, u otro. Dar las credenciales para poder tener acceso.}}

\section{Declaración de contribución de cada integrante}
\textbf{\textit{Describir los aportes de cada integrante al proyecto.}}

\section{Declaración de uso de IA}
\textbf{\textit{Indicar si se emplearon herramientas de inteligencia artificial durante el desarrollo del trabajo. En caso afirmativo, especificar la herramienta utilizada, el propósito de su uso y la forma en que sus aportes fueron revisados, validados e integrados por el grupo. Precisar que la responsabilidad del contenido final recae en los autores.}}

\section{Referencias}
{\color{lightmint}\fontsize{9}{11}\selectfont
Ejemplo:

\textbf{\textit{Victor O.K. Li, “Hints on writing technical papers and making presentations”, IEEE Transactions on Education, vol. 42, no. 2, pp. 134-137, mayo de 1999.}}

\textbf{\textit{Robert M. Woelfle, editor, “New Guide for Better Technical Presentations: Applying Proven Techniques with Modern Tools”, IEEE Professional Communication Society, 1992. Dirección electrónica: www.ieee.org}}
}

\bibliographystyle{IEEEtran}
\bibliography{references}

\subsection*{Nota importante:}
{\color{lightmint}\fontsize{9}{11}\selectfont
\begin{center}
\textbf{Figuras y tablas}
\end{center}
\textbf{\textit{Cada figura debe ser leída y comprendida sin necesidad de recurrir a la descripción a lo largo del texto.}}

\textbf{\textit{Todos los símbolos usados deben ser explicados en la leyenda.}}
}
